\chapter{Product Measures, Fubini, Change of Variables}
\label{ch:b3:product}

One-dimensional Lebesgue theory becomes multi-dimensional
calculus through two theorems. \emph{Tonelli--Fubini} says that
integrals over products are iterated integrals --- slicing is
legitimate, in either order, under hypotheses one can actually
check. The \emph{change of variables formula} transports
integrals along $\mathcal C^1$ diffeomorphisms, with the
Jacobian determinant as the exchange rate for volume; we prove
it completely, starting from the linear case where it explains
what the determinant \emph{is}. Applications cascade: the layer
cake formula, convolution, polar coordinates, the volume of the
$n$-ball --- and, in the weekend problem, Stirling's formula
with an honest error analysis.

\section{Product $\sigma$-algebras and product measures}

\begin{definition}\label{def:b3:product:sigma}
For measurable spaces $(X, \mathcal A)$, $(Y, \mathcal B)$, the
\emph{product $\sigma$-algebra}\index{product
$\sigma$-algebra} $\mathcal A \otimes \mathcal B$ on $X \times
Y$ is generated by the \emph{rectangles} $A \times B$ ($A \in
\mathcal A$, $B \in \mathcal B$) --- a $\pi$-system. For $E
\subseteq X\times Y$ and $x \in X$, the \emph{section} is $E_x
= \{y : (x,y) \in E\}$; for a function $f$ on the product,
$f_x = f(x, \cdot)$.
\end{definition}

\begin{proposition}\label{prop:b3:product:sections}
(a) If $E \in \mathcal A\otimes\mathcal B$, every section $E_x
\in \mathcal B$ (and symmetrically); if $f$ is $\mathcal
A\otimes\mathcal B$-measurable, every $f_x$ is $\mathcal
B$-measurable.
(b) $\mathcal B(\R^m)\otimes\mathcal B(\R^n) = \mathcal
B(\R^{m+n})$.
\end{proposition}

\begin{proof}
(a) Good sets: $\{E : E_x \in \mathcal B\ \forall x\}$ is a
$\sigma$-algebra (sections commute with complements and
countable unions) containing the rectangles. For $f$:
$(f_x)^{-1}(B) = (f^{-1}(B))_x$.
(b) ($\subseteq$) Rectangles of Borel sets: it suffices that
open$\times$open boxes are Borel in $\R^{m+n}$ (they are open)
and that general Borel rectangles are limits --- good sets
again: $\{A : A\times\R^n \in \mathcal B(\R^{m+n})\}$ is a
$\sigma$-algebra containing the opens; intersect two such.
($\supseteq$) Every open set of $\R^{m+n}$ is a countable union
of rational open boxes $U\times V$: contained in the product
$\sigma$-algebra.
\end{proof}

\begin{theorem}[Product measure]\label{thm:b3:product:existence}
Let $(X, \mathcal A, \mu)$ and $(Y, \mathcal B, \nu)$ be
\emph{$\sigma$-finite}. For every $E \in \mathcal
A\otimes\mathcal B$, the function $x \mapsto \nu(E_x)$ is
measurable, and\index{product measure}
\[
(\mu\otimes\nu)(E) = \int_X \nu(E_x)\,\dd\mu(x)
\]
defines the unique measure on $\mathcal A\otimes\mathcal B$
with $(\mu\otimes\nu)(A\times B) = \mu(A)\nu(B)$. It is
$\sigma$-finite, and symmetric: the same measure is obtained by
integrating $x$-sections against $\nu$.
\end{theorem}

\begin{proof}
\emph{Measurability of $x \mapsto \nu(E_x)$.} First let $\nu$
be finite. The class $\mathcal D$ of $E$ for which the map is
measurable contains the rectangles ($\nu((A\times B)_x) =
\nu(B)\mathbf 1_A(x)$) and is a $\lambda$-system: for $E
\subseteq F$ in $\mathcal D$, $\nu((F\setminus E)_x) =
\nu(F_x) - \nu(E_x)$ (finiteness); for $E_n \uparrow E$,
$\nu((E_n)_x) \uparrow \nu(E_x)$ (continuity from below), and
monotone limits of measurable functions are measurable.
Rectangles form a $\pi$-system: Dynkin
(\cref{thm:b3:measure:dynkin}) gives $\mathcal D = \mathcal
A\otimes\mathcal B$. If $\nu$ is $\sigma$-finite, write $Y =
\bigcup Y_k$, $Y_k \uparrow$, $\nu(Y_k) < \infty$: $\nu(E_x) =
\lim_k\nu_k(E_x)$ with $\nu_k = \nu(\cdot\cap Y_k)$ finite.

\emph{Measure.} $\sigma$-additivity of $E \mapsto
\int\nu(E_x)\dd\mu$ follows from
\cref{cor:b3:lebesgue:additivity} (sections of disjoint sets
are disjoint). On rectangles it gives $\mu(A)\nu(B)$.
\emph{Uniqueness}: two candidates agree on the $\pi$-system of
rectangles; $\sigma$-finiteness provides rectangles $X_k\times
Y_k \uparrow X\times Y$ of finite measure:
\cref{thm:b3:measure:uniqueness}. Symmetry: the
other-order construction is also a measure agreeing on
rectangles --- unique, hence the same.
\end{proof}

\begin{definition}\label{def:b3:product:lebesgue}
\emph{Lebesgue measure on $\R^d$}\index{Lebesgue measure on
$\R^d$} is $\lambda_d = \lambda\otimes\cdots\otimes\lambda$
($d$ factors; associativity of the construction is checked on
boxes and propagated by uniqueness). It is the unique Borel
measure giving each box $\prod\intoc{a_i}{b_i}$ its volume
$\prod(b_i - a_i)$; it is translation-invariant (translates
agree on boxes), $\sigma$-finite, and complete after
Carath\'eodory completion --- we write $\lambda_d$ for the
completed measure and integrate accordingly.
\end{definition}

\section{Tonelli and Fubini}

\begin{theorem}[Tonelli]\label{thm:b3:product:tonelli}
$\mu, \nu$ $\sigma$-finite, $f \colon X\times Y \to [0,
+\infty]$ measurable. Then $x \mapsto \int_Y f_x\,\dd\nu$ is
measurable and\index{Tonelli's theorem}
\[
\int_{X\times Y}f\,\dd(\mu\otimes\nu)
= \int_X\Bigl(\int_Y f(x,y)\,\dd\nu(y)\Bigr)\dd\mu(x)
= \int_Y\Bigl(\int_X f(x,y)\,\dd\mu(x)\Bigr)\dd\nu(y).
\]
\end{theorem}

\begin{proof}
The standard machine. For $f = \mathbf 1_E$ this is
\cref{thm:b3:product:existence} (and its symmetric form). By
linearity it holds for simple $f \geq 0$. For general $f \geq
0$: take simple $s_n \nearrow f$
(\cref{thm:b3:lebesgue:approximation}); then $\int_Y(s_n)_x
\dd\nu \nearrow \int_Y f_x\dd\nu$ for each $x$ (MCT in $Y$),
so the left members converge by MCT in $X$, while
$\int s_n\,\dd(\mu\otimes\nu) \nearrow \int f$ by MCT on the
product.
\end{proof}

\begin{theorem}[Fubini]\label{thm:b3:product:fubini}
$\mu, \nu$ $\sigma$-finite, $f \in L^1(\mu\otimes\nu)$. Then
for $\mu$-a.e.\ $x$ the section $f_x$ is $\nu$-integrable, the
a.e.-defined function $x \mapsto \int f_x\dd\nu$ is
integrable, and the two iterated integrals both equal
$\int f\,\dd(\mu\otimes\nu)$.\index{Fubini's theorem}
\end{theorem}

\begin{proof}
Tonelli applied to $\abs f$ shows $\varphi(x) = \int\abs{f_x}
\dd\nu$ has finite integral, hence is finite a.e.: $f_x \in
L^1(\nu)$ for a.e.\ $x$. Split $f = f^+ - f^-$ (real case;
complex by components): Tonelli computes each iterated
integral of $f^\pm$ as $\int f^\pm\dd(\mu\otimes\nu) <
\infty$, and the a.e.-defined difference integrates to the
difference. Symmetrically for the other order.
\end{proof}

\begin{method}\label{met:b3:product:fubini}
To interchange two integrals (or an integral and a sum, or two
sums): if the integrand is \emph{nonnegative}, interchange
freely (Tonelli). Otherwise, first apply Tonelli to $\abs f$
in whichever order is easier to estimate; if the result is
finite, Fubini legitimizes the interchange. Never skip the
$\abs f$ check: \cref{exo:b3:product:4}'s integrand has two
iterated integrals with \emph{different values}.
\end{method}

\begin{proposition}[Layer cake]\label{prop:b3:product:layercake}
For $f \geq 0$ measurable on $(X, \mathcal A, \mu)$
$\sigma$-finite:\index{layer cake formula}
\[
\int_X f\,\dd\mu = \int_0^{+\infty}\mu(\{f > t\})\,\dd t,
\qquad
\int_X f^p\,\dd\mu = p\int_0^{+\infty}t^{p-1}\mu(\{f >
t\})\,\dd t \quad (p \geq 1).
\]
\end{proposition}

\begin{proof}
Apply Tonelli to $\mathbf 1_{\{(x,t) : 0 < t < f(x)\}}$ on $X
\times \intoo0{+\infty}$ (measurable: it is $\{(x,t): f(x) - t
> 0\}\cap\{t > 0\}$, a Borel-type combination of the
measurable $(x,t)\mapsto f(x) - t$): integrating in $t$ first
gives $\int f\,\dd\mu$; in $x$ first, $\int_0^\infty\mu(f >
t)\dd t$. For $f^p$: substitute $t = s^p$ in $\int\mu(f^p > t)
\dd t$, i.e.\ apply the first formula to $f^p$ and change
variables in the one-dimensional integral ($\{f^p > s^p\} =
\{f > s\}$).
\end{proof}

\begin{theorem}[Convolution on $L^1$]\label{thm:b3:product:convolution}
For $f, g \in L^1(\R^d, \lambda_d)$, the integral
\[
(f * g)(x) = \int_{\R^d} f(x - y)\,g(y)\,\dd y
\]
converges absolutely for a.e.\ $x$, defines $f * g \in
L^1(\R^d)$ with $\norm{f*g}_1 \leq \norm f_1\norm g_1$, and
$*$ is commutative and associative.\index{convolution}
\end{theorem}

\begin{proof}
$(x, y) \mapsto f(x-y)g(y)$ is measurable ($ (x,y)\mapsto x -
y$ is continuous; compose and multiply). Tonelli:
\[
\int\!\!\int \abs{f(x-y)}\abs{g(y)}\,\dd y\,\dd x
= \int\abs{g(y)}\Bigl(\int\abs{f(x - y)}\dd x\Bigr)\dd y
= \norm f_1\norm g_1 < \infty
\]
(translation invariance of $\lambda_d$ in the inner integral).
So the double integral is finite; Fubini gives a.e.\ absolute
convergence and the norm bound $\norm{f*g}_1 \leq \norm
f_1\norm g_1$. Commutativity: substitute $y \mapsto x - y$
(translation and reflection invariance --- reflection
invariance holds on boxes, hence everywhere by uniqueness).
Associativity: Tonelli--Fubini on a triple integral.
\end{proof}

\section{Change of variables}

\begin{theorem}[Linear change of variables]
\label{thm:b3:product:linearchange}
For $T \in GL_d(\R)$ and $A \in \mathcal B(\R^d)$:
$\lambda_d(T(A)) = \abs{\det T}\,\lambda_d(A)$; consequently
$\int f(y)\dd y = \abs{\det T}\int f(Tx)\,\dd x$ for $f \geq 0$
or integrable.\index{linear change of variables}
\end{theorem}

\begin{proof}
The measure $\mu_T(A) = \lambda_d(T(A))$ is a Borel measure
(homeomorphisms preserve Borel sets,
\cref{pb:b3:measure:1}), translation-invariant
($T(A + x) = T(A) + Tx$), finite on the unit box: by the
characterization of Lebesgue measure
(\cref{exo:b3:measure:6}, whose proof works verbatim in
$\R^d$ with dyadic cubes), $\mu_T = c(T)\lambda_d$ with $c(T)
= \lambda_d(T(\intco01^d))$. The map $T \mapsto c(T)$ is
multiplicative ($c(ST) = c(S)c(T)$, by composing), so it
suffices to compute $c$ on generators of $GL_d$: elementary
matrices. Diagonal $\operatorname{diag}(a, 1, \dots, 1)$:
maps the unit cube to a box of volume $\abs a$: $c = \abs a =
\abs\det$. Transposition of coordinates: permutes the cube: $c
= 1 = \abs\det$. Transvection $T(x) = x + \alpha
x_2e_1$: the image of the unit cube is a sheared prism; by
Tonelli its measure is $\int\lambda_1(\text{section})\dd x_2
\cdots \dd x_d = 1$, each $x_1$-section being an interval of
length $1$: $c = 1 = \abs{\det}$. Every invertible matrix is a
product of these (Gaussian elimination), and both $c$ and
$\abs\det$ are multiplicative: $c(T) = \abs{\det T}$. The
integral formula follows by the standard machine (indicators,
simple, MCT).
\end{proof}

\begin{theorem}[Change of variables]\label{thm:b3:product:changeofvar}
Let $U, V \subseteq \R^d$ be open and $\Phi \colon U \to V$ a
$\mathcal C^1$ diffeomorphism. For every measurable $f \colon V
\to [0, +\infty]$ (or $f \in L^1(V)$):\index{change of
variables formula}
\[
\int_V f(y)\,\dd y = \int_U
f\bigl(\Phi(x)\bigr)\,\abs{\det D\Phi(x)}\,\dd x .
\]
\end{theorem}

\begin{proof}
Write $J(x) = \abs{\det D\Phi(x)}$. The heart of the proof is
the inequality
\[
\lambda_d\bigl(\Phi(A)\bigr) \leq \int_A J\,\dd\lambda_d
\qquad\text{for every Borel } A \subseteq U;
\tag{$*$}
\]
Step 4 below upgrades $(*)$ --- applied to both $\Phi$ and
$\Phi^{-1}$ --- to the equality of the theorem. Note that
$\Phi^{-1}$ is itself a $\mathcal C^1$ diffeomorphism with
Jacobian $\abs{\det D\Phi^{-1}(y)} = J(\Phi^{-1}y)^{-1}$
(chain rule on $\Phi\circ\Phi^{-1} = \mathrm{id}$).

\emph{Step 1: $(*)$ for cubes with a distortion factor.} Fix a
closed cube $Q \subseteq U$ of center $x_0$ and side $2r$
(sup-norm ball). Claim: for every $\varepsilon > 0$, if $\Phi$
is differentiable on $Q$ with
$\norm{D\Phi(x) - D\Phi(x_0)} \leq \varepsilon$ on $Q$
(operator norm for the sup-norm), then
\[
\Phi(Q) \subseteq \Phi(x_0) + D\Phi(x_0)\Bigl(\,\bigl(1 +
\varepsilon\norm{D\Phi(x_0)^{-1}}\bigr)\,(Q - x_0)\Bigr),
\]
because for $x \in Q$, the mean value inequality applied to
$\Phi(x) - \Phi(x_0) - D\Phi(x_0)(x - x_0)$ gives
$\norm{\Phi(x) - \Phi(x_0) - D\Phi(x_0)(x-x_0)}_\infty \leq
\varepsilon\norm{x - x_0}_\infty \leq \varepsilon r$, and
$D\Phi(x_0)^{-1}$ pulls this defect into an
$\varepsilon\norm{D\Phi(x_0)^{-1}}\,r$-enlargement of the cube.
By \cref{thm:b3:product:linearchange},
\[
\lambda_d(\Phi(Q)) \leq \abs{\det D\Phi(x_0)}\,
\bigl(1 + \varepsilon\,C\bigr)^{d}\,\lambda_d(Q),
\qquad C = \sup_{Q}\norm{D\Phi(\cdot)^{-1}} .
\]

\emph{Step 2: $(*)$ for compact cubes, by subdivision.} Let $Q
\subseteq U$ be a compact cube and $\varepsilon > 0$. On $Q$,
$D\Phi$ is uniformly continuous and $\norm{D\Phi^{-1}}$
bounded (compactness); subdivide $Q$ into $2^{kd}$ subcubes
$Q_i$ small enough that the oscillation of $D\Phi$ on each is
$\leq \varepsilon$. Step 1 on each subcube (center $x_i$):
\[
\lambda_d(\Phi(Q)) \leq \sum_i\lambda_d(\Phi(Q_i))
\leq (1 + C\varepsilon)^d \sum_i \abs{\det
D\Phi(x_i)}\,\lambda_d(Q_i)
\leq (1 + C\varepsilon)^d\Bigl(\int_Q J + \varepsilon'\Bigr),
\]
the last step because $\sum_i\abs{\det D\Phi(x_i)}\mathbf
1_{Q_i} \to J$ uniformly on $Q$ (continuity of $\det D\Phi$)
--- Riemann-sum comparison. Let $\varepsilon \to 0$: $(*)$
holds for compact cubes.

\emph{Step 3: $(*)$ for all Borel $A$.} The set function $A
\mapsto \lambda_d(\Phi(A))$, on Borel subsets of $U$, is a
measure ($\Phi$ is a bijection onto $V$ preserving Borel sets
and countable disjointness), and so is $A \mapsto \int_AJ$.
Every open subset of $U$ is a countable union of
almost-disjoint dyadic compact cubes (standard dyadic
decomposition: take maximal dyadic cubes contained in the open
set), and both measures are additive across them (boundaries
of cubes are $\lambda_d$-null, and their $\Phi$-images are
null by Step 2 applied to thin cube-coverings of the faces):
$(*)$ passes from cubes to open sets. General Borel $A$:
exhaust $U$ by compacts $K_m \uparrow U$ with $K_m \subseteq
\mathring K_{m+1}$, and fix $m$; on $\mathring K_{m+1}$, $J$
is bounded by some $M_m$. By outer regularity of $\lambda_d$
(proof as in \cref{thm:b3:measure:regularity}, with boxes),
choose open sets $O_n$ with $A \cap K_m \subseteq O_n
\subseteq \mathring K_{m+1}$ and $\lambda_d\bigl(O_n \setminus
(A\cap K_m)\bigr) \to 0$. Then
\[
\lambda_d\bigl(\Phi(A\cap K_m)\bigr) \leq
\lambda_d\bigl(\Phi(O_n)\bigr) \leq \int_{O_n}J
\leq \int_{A\cap K_m}J + M_m\,\lambda_d\bigl(O_n\setminus(A\cap
K_m)\bigr) \xrightarrow[n\to\infty]{} \int_{A\cap K_m}J .
\]
Let $m \to \infty$: continuity from below on the left, MCT on
the right. This establishes $(*)$.

\emph{Step 4: equality and the integral formula.} First
extend $(*)$ from sets to integrals: for every measurable $g
\geq 0$ on $V$,
\[
\int_V g(y)\,\dd y \leq \int_U g(\Phi(x))\,J(x)\,\dd x .
\tag{$**$}
\]
Indeed, for $g = \mathbf 1_B$ this is $(*)$ with $A =
\Phi^{-1}(B)$; linearity extends it to simple $g$, and MCT to
all $g \geq 0$ (the standard machine). Now apply $(**)$ twice:
first to $g$, then --- for the diffeomorphism $\Phi^{-1}$ ---
to the function $x \mapsto g(\Phi(x))J(x)$:
\[
\int_Vg \leq \int_U g(\Phi(x))J(x)\dd x
\leq \int_V g(y)\,J(\Phi^{-1}y)\,\abs{\det
D\Phi^{-1}(y)}\,\dd y = \int_V g ,
\]
since $J(\Phi^{-1}y)\abs{\det D\Phi^{-1}(y)} = \abs{\det\bigl(
D\Phi(\Phi^{-1}y)\,D\Phi^{-1}(y)\bigr)} = 1$ (chain rule on
$\Phi\circ\Phi^{-1} = \mathrm{id}$). All inequalities are
equalities: the formula holds for $g \geq 0$, and for $L^1$
functions by decomposition.
\end{proof}

\begin{example}[Polar coordinates; the Gaussian again]
\label{ex:b3:product:polar}
$\Phi(r, \theta) = (r\cos\theta, r\sin\theta)$ is a $\mathcal
C^1$ diffeomorphism from $\intoo0{+\infty}\times\intoo0{2\pi}$
onto $\R^2$ minus a half-line (null set), with $\det D\Phi =
r$:\index{polar coordinates}
\[
\int_{\R^2}f(x, y)\,\dd x\,\dd y
= \int_0^{2\pi}\!\!\int_0^{+\infty}
f(r\cos\theta, r\sin\theta)\,r\,\dd r\,\dd\theta .
\]
For $f = \eu^{-x^2-y^2}$, Tonelli and this formula give
\[
G^2 = \Bigl(\int_\R \eu^{-x^2}\dd x\Bigr)^2
= \int_{\R^2}\eu^{-x^2-y^2}
= 2\pi\int_0^\infty r\eu^{-r^2}\dd r = \pi:
\]
the classical two-line proof of $G = \sqrt\pi$, now fully
justified (compare the parameter proof of
\cref{pb:b3:lebesgue:1}).
\end{example}

\begin{theorem}[Volume of the unit ball]
\label{thm:b3:product:ballvolume}
Let $v_d = \lambda_d(B(0,1))$ in $\R^d$. Then\index{volume of
the unit ball}
\[
v_d = \frac{\pi^{d/2}}{\Gamma\bigl(\frac d2 + 1\bigr)} :
\qquad
v_1 = 2,\quad v_2 = \pi,\quad v_3 = \tfrac{4\pi}3,\quad
v_4 = \tfrac{\pi^2}2,\ \dots
\]
\end{theorem}

\begin{proof}
Compute $I = \int_{\R^d}\eu^{-\norm x_2^2}\dd\lambda_d$ twice.
By Tonelli it factors: $I = G^d = \pi^{d/2}$. By the layer cake
formula (\cref{prop:b3:product:layercake}) with $f = \eu^{-
\norm x^2}$, whose level sets are balls: $\{f > t\} =
B\bigl(0, \sqrt{-\ln t}\bigr)$ for $0 < t < 1$, of measure
$v_d(-\ln t)^{d/2}$ (dilation by $\rho$ scales $\lambda_d$ by
$\rho^d$: \cref{thm:b3:product:linearchange}), so
\[
I = \int_0^1 v_d\,(-\ln t)^{d/2}\,\dd t
\overset{t = \eu^{-s}}{=}
v_d\int_0^\infty s^{d/2}\eu^{-s}\,\dd s
= v_d\,\Gamma\Bigl(\frac d2 + 1\Bigr).
\]
Equate. (The values: $\Gamma(\frac32) = \frac{\sqrt\pi}2$,
$\Gamma(2) = 1$, etc.) Note $v_d \to 0$ as $d \to \infty$ ---
the weekend problem quantifies how fast, via Stirling.
\end{proof}

\section{Exercises}

\begin{exercise}[$\star$]\label{exo:b3:product:1}
Let $\mu$ be counting measure on $(\intcc01, \mathcal
B(\intcc01))$ (not $\sigma$-finite) and $\lambda$ Lebesgue
measure, and let $\Delta = \{(x,x)\}$ be the diagonal in
$\intcc01^2$. Show that $\Delta$ is measurable, and compute the
two iterated integrals of $\mathbf 1_\Delta$ against
$\lambda$ and $\mu$: they differ. Which hypothesis of
\cref{thm:b3:product:tonelli} fails?
\end{exercise}

\begin{exercise}[$\star$]\label{exo:b3:product:2}
Justify the interchange and re-derive Dirichlet's integral: for
$A > 0$,
\[
\int_0^A\frac{\sin x}x\,\dd x
= \int_0^A\!\!\int_0^{+\infty}\eu^{-xy}\sin x\,\dd y\,\dd x
= \int_0^{+\infty}\!\!\int_0^A \eu^{-xy}\sin x\,\dd x\,\dd y,
\]
compute the inner integral in closed form, and let $A \to
+\infty$ (dominate the $y$-integral) to get
$\int_0^\infty\frac{\sin x}x\dd x = \frac\pi2$.
\end{exercise}

\begin{exercise}[$\star\star$]\label{exo:b3:product:3}
(a) Prove that for $f \geq 0$ measurable and $\mu$ finite:
$\sum_{n\geq1}\mu(\{f \geq n\}) \leq \int f\,\dd\mu \leq
\mu(X) + \sum_{n\geq1}\mu(\{f\geq n\})$: integrability is
summability of the tail measures.
(b) Deduce that $f \in L^1(\mu)$ ($\mu$ finite) iff
$\sum_n\mu(\abs f \geq n) < \infty$.
\end{exercise}

\begin{exercise}[$\star\star$]\label{exo:b3:product:4}
For $f(x, y) = \dfrac{x^2 - y^2}{(x^2 + y^2)^2}$ on
$\intoo01^2$, show
\[
\int_0^1\!\!\int_0^1 f\,\dd y\,\dd x = \frac\pi4,
\qquad
\int_0^1\!\!\int_0^1 f\,\dd x\,\dd y = -\frac\pi4
\]
\emph{(note $f = \partial_y\bigl(\frac{y}{x^2+y^2}\bigr)$)},
and verify directly that $\int\!\int\abs f = +\infty$: Fubini's
integrability hypothesis is not decorative.
\end{exercise}

\begin{exercise}[$\star\star$]\label{exo:b3:product:5}
(a) Compute $\mathbf 1_{\intcc01} * \mathbf 1_{\intcc01}$
explicitly (a tent function), and $(\mathbf 1 * \mathbf 1 *
\mathbf 1)$'s general shape.
(b) Show $\operatorname{supp}(f * g) \subseteq
\overline{\operatorname{supp}f + \operatorname{supp}g}$.
(c) Show that if $f \in L^1$ and $g$ is bounded and continuous,
$f * g$ is continuous. \emph{(DCT via continuity of
translation on the bounded $g$.)}
\end{exercise}

\begin{exercise}[$\star\star$]\label{exo:b3:product:6}
(a) Show that the simplex $\Delta_d = \{x \in \intco0\infty^d
: x_1 + \dots + x_d \leq 1\}$ has volume $\frac1{d!}$
\emph{(induction and Fubini)}.
(b) Recover $v_2 = \pi$, $v_3 = \frac{4\pi}3$ from
\cref{thm:b3:product:ballvolume}, and show
$\lambda_d(\text{ellipsoid with semi-axes } a_i) =
v_d\prod a_i$.
\end{exercise}

\begin{exercise}[$\star\star$]\label{exo:b3:product:7}
For which $s > 0$ are the following finite? Justify with polar
coordinates:
\[
\int_{B(0,1)\subseteq\R^2}\frac{\dd x\,\dd y}{(x^2 +
y^2)^{s}},
\qquad
\int_{\R^2\setminus B(0,1)}\frac{\dd x\,\dd y}{(x^2 +
y^2)^{s}} .
\]
Generalize to $\R^d$ (the thresholds $s < d/2$ and $s > d/2$).
\end{exercise}

\begin{exercise}[$\star\star\star$]\label{exo:b3:product:8}
(Beta--Gamma) For $p, q > 0$, let $B(p, q) = \int_0^1t^{p-1}(1
- t)^{q-1}\dd t$. Starting from $\Gamma(p)\Gamma(q)$ as a
double integral, substitute $(x, y) = (uv,\, u(1 - v))$ (a
diffeomorphism of the open quadrant onto
$\intoo0\infty\times\intoo01$; compute its Jacobian $= u$) and
conclude\index{Beta function}
\[
B(p, q) = \frac{\Gamma(p)\,\Gamma(q)}{\Gamma(p + q)} .
\]
Deduce $\int_0^{\pi/2}\sin^{2p-1}\theta\cos^{2q-1}\theta\,
\dd\theta = \frac12B(p,q)$ and the value of the Wallis
integrals $W_n = \int_0^{\pi/2}\sin^n$.
\end{exercise}

\begin{exercise}[$\star\star$]\label{exo:b3:product:9}
(Transfer formula) Let $T \colon (X, \mathcal A, \mu) \to (Y,
\mathcal B)$ be measurable and $T_*\mu(B) = \mu(T^{-1}(B))$
the \emph{pushforward measure}\index{pushforward measure}.
Show that for every measurable $g \geq 0$ on $Y$:
\[
\int_Y g\,\dd(T_*\mu) = \int_X g\circ T\,\dd\mu
\]
(standard machine). Then compare with
\cref{thm:b3:product:linearchange}: what extra information does
the change of variables formula carry that the abstract
transfer formula does not? \emph{(The transfer formula never
identifies $T_*\mu$; the change of variables theorem computes
$T_*\lambda_d$ explicitly as a density measure.)}
\end{exercise}

\begin{exercise}[$\star\star\star$]\label{exo:b3:product:10}
(Gaussian moments) Using polar coordinates and Fubini, compute
for the standard Gaussian weight on $\R^d$:
\[
\int_{\R^d}\norm x_2^2\;\eu^{-\norm x_2^2}\,\dd x
\qquad\text{and}\qquad
\int_{\R^d}x_1^2\,\eu^{-\norm x^2_2}\,\dd x,
\]
check the consistency ($\norm x^2 = \sum x_i^2$), and deduce
the second moment of the measure
$\pi^{-d/2}\eu^{-\norm x^2}\dd x$.
\end{exercise}

\begin{exercise}[$\star\star$]\label{exo:b3:product:11}
(Graph and hypograph) Let $f \colon \R^d \to \intco0\infty$
be measurable.
(a) Show that the \emph{hypograph} $H = \{(x, y) \in
\R^d\times\R : 0 < y < f(x)\}$ is measurable in
$\R^{d+1}$ with
\[
\lambda_{d+1}(H) = \int_{\R^d}f\,\dd\lambda_d :
\]
``the integral is the area under the graph'', at last a
theorem. \emph{(Sections; Tonelli.)}
(b) Show that the graph $\{(x, f(x)) : x \in \R^d\}$ is a
null set of $\R^{d+1}$.
(c) Deduce a two-line proof that the sphere $S^{d-1}$ is
Lebesgue-null in $\R^d$.
\end{exercise}

\begin{exercise}[$\star\star$]\label{exo:b3:product:12}
(A famous double integral) Using the geometric series and
Tonelli on $\intoo01^2$, prove
\[
\int_0^1\!\!\int_0^1\frac{\dd x\,\dd y}{1 - xy}
= \sum_{n\geq1}\frac1{n^2} = \zeta(2),
\qquad
\int_0^1\!\!\int_0^1\frac{\dd x\,\dd y}{1 + xy}
= \sum_{n\geq1}\frac{(-1)^{n-1}}{n^2} = \frac{\zeta(2)}2 .
\]
(The second series identity: split even and odd indices.)
With $\zeta(2) = \frac{\pi^2}6$ (\cref{ch:b3:spectral}), two
innocent-looking integrals evaluate to $\frac{\pi^2}6$ and
$\frac{\pi^2}{12}$; where exactly does Tonelli's positivity
hypothesis do its work?
\end{exercise}

\section{Problem: Stirling's formula}

\begin{problem}[{Weekend problem --- $n! \sim
\sqrt{2\pi n}\,(n/\eu)^n$, by dominated convergence}]
\label{pb:b3:product:1}
Stirling's formula governs every asymptotic count in this book
--- ball volumes, binomial coefficients, the central limit
theorem's local form. We prove it from the $\Gamma$ integral
(\cref{ex:b3:lebesgue:gamma}) with the \emph{Laplace method},
in its cleanest dominated-convergence form, then collect
dividends.\index{Stirling's formula}

\medskip
\noindent\textbf{Part I --- The formula.} For $t > 0$,
$\Gamma(t + 1) = \int_0^\infty x^{t}\eu^{-x}\dd x$.
\begin{enumerate}
  \item Substitute $x = t + \sqrt t\,u$ and show
        \[
        \frac{\Gamma(t+1)}{t^{t}\eu^{-t}\sqrt t}
        = \int_{-\sqrt t}^{+\infty}
        \exp\Bigl(t\ln\Bigl(1 + \frac u{\sqrt
        t}\Bigr) - \sqrt t\,u\Bigr)\,\dd u
        \;=\;\int_\R g_t(u)\,\dd u,
        \]
        where $g_t(u) = \exp\bigl(t\ln(1 + u/\sqrt t) -
        \sqrt t\,u\bigr)\mathbf 1_{u > -\sqrt t}$.
  \item Show the pointwise limit: for every fixed $u$, $g_t(u)
        \to \eu^{-u^2/2}$ as $t \to +\infty$ \emph{(expand
        $\ln(1 + h)$ to second order)}.
  \item Domination. Let $\varphi(h) = \ln(1 + h) - h$, so that
        $g_t(u) = \exp\bigl(t\,\varphi(u/\sqrt t)\bigr)$ for
        $u > -\sqrt t$. Prove the two bounds
        \[
        \varphi(h) \leq -\frac{h^2}4 \quad (-1 < h \leq 1),
        \qquad
        \varphi(h) \leq -c\,h \quad (h \geq 1),\ \ c = 1 -
        \ln 2 > 0
        \]
        \emph{(study $\varphi(h) + \frac{h^2}4$ and $\varphi(h)
        + ch$: compute the derivatives and check the sign on
        each range)}. Deduce, for $t \geq 1$:
        \[
        g_t(u) \leq \eu^{-u^2/4}\ \ (\abs u \leq \sqrt t),
        \qquad
        g_t(u) \leq \eu^{-cu}\ \ (u \geq \sqrt t),
        \]
        so that $g_t(u) \leq \eu^{-u^2/4} +
        \eu^{-cu}\,\mathbf 1_{u > 0}$: an integrable dominator
        independent of $t \geq 1$.
  \item Conclude with the dominated convergence theorem and the
        Gaussian integral (\cref{ex:b3:product:polar}):
        \[
        \Gamma(t + 1) \;\sim\;
        \sqrt{2\pi t}\;\Bigl(\frac t\eu\Bigr)^{t}
        \qquad (t \to +\infty),
        \]
        and in particular $n! \sim \sqrt{2\pi
        n}\,(n/\eu)^n$.
\end{enumerate}

\medskip
\noindent\textbf{Part II --- Dividends.}
\begin{enumerate}[resume]
  \item (Wallis) From \cref{exo:b3:product:8}, $W_{2n} =
        \frac\pi2\binom{2n}n4^{-n}$-type formulas: derive
        $\binom{2n}{n} \sim \frac{4^n}{\sqrt{\pi n}}$ from
        Stirling, and check it against the recursion $W_{n} =
        \frac{n-1}nW_{n-2}$.
  \item (Ball volumes collapse) Show
        \[
        v_d = \frac{\pi^{d/2}}{\Gamma(\frac d2 + 1)}
        \;\sim\; \frac{1}{\sqrt{\pi d}}
        \Bigl(\frac{2\pi\eu}{d}\Bigr)^{d/2},
        \]
        so $v_d \to 0$ faster than any geometric sequence;
        find the dimension maximizing $v_d$ (numerically: $d =
        5$).
  \item (Concentration of the binomial --- a preview of
        \cref{ch:b3:clt}) Using Stirling, show the local
        estimate, for $k = n/2 + s\sqrt n/2$ with $s$ fixed
        and $n$ even:
        \[
        2^{-n}\binom{n}{k} \;\sim\;
        \sqrt{\frac{2}{\pi n}}\;\eu^{-s^2/2},
        \]
        the discrete Gaussian profile: de Moivre--Laplace in
        embryo.
  \item Where exactly did the proof of Part I use: (i) MCT or
        DCT; (ii) the Gaussian integral; (iii) the invariance
        properties of Lebesgue measure? One sentence each.
\end{enumerate}

\medskip
\noindent\textbf{Part III --- The error term: Stirling with
bars.} Set $d_n = \ln n! - \bigl(n + \tfrac12\bigr)\ln n + n -
\ln\sqrt{2\pi}$, so that Part I says $d_n \to 0$.
\begin{enumerate}[resume]
  \item Show $d_n - d_{n+1} = \bigl(n +
        \tfrac12\bigr)\ln\bigl(1 + \tfrac1n\bigr) - 1$.
  \item With $t = \frac1{2n+1}$, verify $\frac{n+1}n =
        \frac{1+t}{1-t}$ and expand:
        \[
        d_n - d_{n+1} = \frac{t^2}3 + \frac{t^4}5 +
        \frac{t^6}7 + \cdots,
        \]
        and deduce the two-sided bounds
        \[
        \frac1{3(2n+1)^2} \;<\; d_n - d_{n+1} \;<\;
        \frac1{12n} - \frac1{12(n+1)} .
        \]
  \item Telescope (using $d_m \to 0$) and check the pleasant
        algebraic identity $\frac1{3(2m+1)^2} >
        \frac1{12m+1} - \frac1{12(m+1)+1}$ for $m \geq 1$, to
        obtain the classical bracketing
        \[
        \sqrt{2\pi n}\Bigl(\frac n\eu\Bigr)^n
        \eu^{1/(12n+1)} \;<\; n! \;<\;
        \sqrt{2\pi n}\Bigl(\frac n\eu\Bigr)^n\eu^{1/(12n)} .
        \]
  \item Two consequences: (a) the relative error of Stirling's
        formula is $< 10^{-6}$ as soon as $n \geq 83\,334$;
        (b) estimate $100!$ to four significant digits by hand
        from the bracketing ($100! \approx 9.3326\cdot
        10^{157}$), and marvel briefly at the precision of an
        asymptotic formula at a very finite $n$.
\end{enumerate}

\medskip
\noindent\textbf{Part IV --- The Wallis route: Stirling
without the Gaussian.} Historically the constant $\sqrt{2\pi}$
came from Wallis, not from Gauss; this part re-proves Stirling
independently of Parts I--II, and thereby re-proves the
Gaussian integral. Let $W_n = \int_0^{\pi/2}\sin^n\theta\,
\dd\theta$.
\begin{enumerate}[resume]
  \item Establish $W_n = \frac{n-1}nW_{n-2}$ (integrate by
        parts), the closed forms
        \[
        W_{2n} = \frac\pi2\binom{2n}n4^{-n}, \qquad
        W_{2n+1} = \frac{4^n}{(2n+1)\binom{2n}n},
        \]
        and the identity $W_nW_{n-1} = \frac\pi{2n}$.
  \item From the monotonicity of $(W_n)$ deduce
        $W_{2n}/W_{2n+1} \to 1$, then
        \[
        W_{2n} \sim \frac12\sqrt{\frac\pi n}
        \qquad\text{and}\qquad
        \binom{2n}n4^{-n}\sqrt n \longrightarrow
        \frac1{\sqrt\pi} :
        \]
        Wallis' theorem, obtained without Stirling.
  \item Show, by the telescoping of Part III alone (no value
        of the constant needed), that $e_n = \ln n! - (n +
        \frac12)\ln n + n$ converges to some limit $\ell$;
        equivalently $n! \sim K\,n^{n+1/2}\eu^{-n}$ with $K =
        \eu^\ell > 0$ not yet identified.
  \item Insert this asymptotic into $\binom{2n}n4^{-n}\sqrt n$
        and identify, using question 14, the only possible
        value: $K = \sqrt{2\pi}$. Assemble the logic: Parts
        III--IV together give a complete second proof of
        Stirling --- and hence, running Part I's substitution
        backwards, an independent evaluation of
        $\int_\R\eu^{-u^2/2}\dd u = \sqrt{2\pi}$. Two pillars,
        either of which supports the other.
\end{enumerate}

\medskip
\noindent\textbf{Part V --- Last dividends.}
\begin{enumerate}[resume]
  \item (The full local profile) For integers $\abs j \leq
        K\sqrt n$ ($K$ fixed), show
        \[
        \frac{\binom{2n}{n+j}}{\binom{2n}{n}} =
        \prod_{i=1}^{\abs j}\frac{n - i + 1}{n + i}
        = \exp\Bigl(-\frac{j^2}n +
        O\Bigl(\frac1{\sqrt n}\Bigr)\Bigr),
        \]
        uniformly in $j$ \emph{(take logarithms and use $\ln
        \frac{1-x}{1+y} = -(x + y) + O(x^2 + y^2)$)}. This is
        the two-sided version of question 7 and the exact
        estimate quoted in \cref{ch:b3:clt}'s weekend problem.
  \item (A Poisson preview) Show with Stirling that
        $\eu^{-n}\dfrac{n^n}{n!} \sim \dfrac1{\sqrt{2\pi n}}$:
        the mode of a Poisson law of large mean $n$ carries
        mass $\approx (2\pi n)^{-1/2}$, exactly as the central
        limit theorem will predict.
  \item (Gamma ratios) For $a \in \intoo01$, prove
        $\dfrac{\Gamma(n + a)}{\Gamma(n)\,n^a} \to 1$ using
        the log-convexity slope bounds of
        \cref{pb:b3:lebesgue:1} (question 14 there), and
        extend to every real $a > 0$ by the functional
        equation. (This is what ``$\Gamma(t+1) \sim$
        Stirling'' means between the integers.)
  \item (Balls, encore) From $v_d = \pi^{d/2}/\Gamma(\frac d2
        + 1)$: tabulate $v_1, \dots, v_7$ exactly, verify
        unimodality via $\frac{v_d}{v_{d-2}} = \frac{2\pi}d$
        (increasing while $d < 2\pi$, decreasing after), and
        prove the striking generating identity
        \[
        \sum_{k\geq0}v_{2k}\,x^{2k} = \eu^{\pi x^2} :
        \]
        all even-dimensional unit-ball volumes packed into one
        exponential.
  \item (Entropy asymptotics) For fixed $\alpha \in \intoo01$
        with $\alpha n \in \N$, deduce from Stirling
        \[
        \binom{n}{\alpha n} \;\sim\;
        \frac{\eu^{n\,H(\alpha)}}
        {\sqrt{2\pi\,\alpha(1-\alpha)\,n}},
        \qquad
        H(\alpha) = -\alpha\ln\alpha -
        (1-\alpha)\ln(1-\alpha) :
        \]
        the exponential growth rate of binomial coefficients
        is the \emph{entropy} $H$ --- check that $\alpha =
        \frac12$ recovers question 5, and that $H(\alpha) <
        \ln2$ for $\alpha \neq \frac12$ (so off-center
        binomials are exponentially negligible in $2^n$).
  \item (Surface areas) The area of the unit sphere $S^{d-1}$
        is $s_{d-1} = d\,v_d$ (proved as
        \cref{exo:b3:forms:6} in the differential-forms
        chapter; here, take it as the definition). Tabulate
        $s_0, \dots, s_6$, locate the maximal one ($d - 1 =
        6$, $s_6 = \frac{16\pi^3}{15} \approx 33.07$), and
        show $s_{d-1} \to 0$ super-geometrically as well ---
        high-dimensional spheres are, by every Euclidean
        yardstick, vanishingly small.
  \item (The first correction term) Deduce from the
        bracketing of question 11 that $d_n = \frac1{12n} +
        O\bigl(\frac1{n^2}\bigr)$, hence
        \[
        n! = \sqrt{2\pi n}\,\Bigl(\frac
        n\eu\Bigr)^{n}\Bigl(1 + \frac1{12n} +
        O\Bigl(\frac1{n^2}\Bigr)\Bigr).
        \]
        Verify at $n = 10$: the bare formula gives
        $3\,598\,696$ (relative error $8.3\cdot10^{-3}$),
        the corrected one $3\,628\,685$ against $10! =
        3\,628\,800$ (relative error $3.2\cdot10^{-5}$) ---
        one term of the series buys two and a half digits.
  \item (The median of $\Gamma$) Show that
        \[
        \frac{1}{\Gamma(t+1)}
        \int_0^{t} x^{t}\eu^{-x}\,\dd x
        \;\longrightarrow\; \frac12
        \qquad (t \to +\infty) :
        \]
        asymptotically, exactly half of the mass of the
        $\Gamma$ integrand sits below its mode $x = t$.
        \emph{(Run Part I's substitution on the truncated
        integral; the dominator of question 3 is already in
        place.)}
  \item (Entropy, non-asymptotically) For $\alpha \in
        \intoc0{\frac12}$ prove the bound, valid for
        \emph{every} $n \geq 1$:
        \[
        \sum_{k=0}^{\lfloor\alpha n\rfloor}\binom nk
        \;\leq\; \eu^{n\,H(\alpha)} ,
        \]
        by comparing the sum with $\sum_k\binom
        nk\lambda^{k-\alpha n}$ for the tilt $\lambda =
        \frac{\alpha}{1-\alpha} \leq 1$. Check that this
        choice of $\lambda$ is optimal, and reconcile with
        question 21: the exponential rate $H(\alpha)$ of the
        asymptotic statement is attained by a one-line
        inequality with no asymptotics at all.
\end{enumerate}
\end{problem}
